\section{Essential heights on every positive residence scale}
\label{sec:v85-bands}
The residence argument need not stop at the long segments. On a fixed
residence band the entire segment lies near any roof that it contains.
A wider exit kernel therefore persists across the whole segment. This
controls all endpoint types on that band, including the critical and
two-nongrazing-end types left in the preceding reduction.

For a nongrazing outgoing face write $b=F(y)$, $l=\tau^-(y)$ and
$a=b-l$. In the parametrization of Theorem~\ref{thm:v85-segment-map},
define
\begin{align}
 k_r(t;y)&=\left(1-\frac{b-t}{r}\right)_+
                                    \one_{\{a<t<b\}},\notag\\
 \mathcal B_s(t)&=\int_{\Gamma_{\rm ng}^+}
          \one_{\{s<l\le2s\}}\one_{\{a<t<b\}}\,d\zeta(y).
                                             \label{eq:v85-band-source}
\end{align}
Thus $\overline{\mathcal F}_r=\int k_r\,d\zeta$. The band source
uses the full occupation indicator of the same maximal segment, not a
renormalized boundary measure. A superscript $[b_0]$ will restrict an
outgoing face to $z(b-)<b_0$.

\begin{theorem}[Complete residence-band heights]
\label{thm:v85-band-height}
For each fixed sufficiently small $s>0$, almost everywhere in $t$,
\begin{equation}\label{eq:v85-band-persistence}
 \mathcal B_s(t)\le\frac2s\int_{t-2s}^{t+2s}
                                     \overline{\mathcal F}_{4s}(v)\,dv.
\end{equation}
Consequently, with $\sigma=1/32$,
\begin{equation}\label{eq:v85-band-height}
 \limsup_m\sup_{R,n,k}m^2\|\mathcal B_s\|_\infty
                                                   \le C_\eta s^\sigma.
\end{equation}
All nongrazing outgoing faces in the band are included, irrespective
of their incoming endpoint or outgoing gradient. More precisely,
for every fixed small $b_0>0$,
\begin{align}
 \limsup_m\sup_{R,n,k}m^2\|\mathcal B_s^{[b_0]}\|_\infty
 \le{}& Cb_0^\alpha
       +C_\eta s^{(1+\eta)/2}b_0^{\alpha-\eta}.
                                             \label{eq:v85-band-critical}
\end{align}
It also satisfies \eqref{eq:v85-band-height}. In particular the
complete band with $z(b-)^2<8\Lambda_\circ s$ has essential height
$O_\eta(s^{129/256})$, not merely local mass of that order.
\end{theorem}
\begin{proof}
If $a<t<b$ and $s<l\le2s$, then $(a,b)\subset(t-2s,t+2s)$.
For every $v\in(a,b)$,
$k_{4s}(v;y)=1-(b-v)/(4s)\ge1/2$.
Thus the integral of this kernel over that window is at least
$l/2>s/2$. Multiply by $2/s$ and integrate over the same faces.
Tonelli proves \eqref{eq:v85-band-persistence}; removing the band
restriction only on the upper side preserves the actual source.
Theorem~\ref{thm:v83-surface-local-mass}, at width $4s$ and interval
length $4s$, now gives
$C s^{1/32}+C_\eta s^{(1+\eta)/2}$ after division by $s$.
This proves \eqref{eq:v85-band-height}. No pointwise density bound
has been inferred from mass without the preceding persistence.

Retain the outgoing restriction $z(b-)<b_0$ on both sides instead.
Radial monotonicity keeps the same restriction at every interior
point of these segments. The positive-variation surface comparison
of Theorem~\ref{thm:v83-two-sided-flux}, restricted to these faces,
bounds the surface by the small-gradient starting source plus
$C_\eta\mathsf P_{\eta,4s}v^{\eta,b_0}$. Its mass on a window
of length $4s$ is at most
$Cs b_0^\alpha+C_\eta s^{1+(1+\eta)/2}b_0^{\alpha-\eta}$
in the prescribed ordered norm. The starting-source bound follows
from \eqref{eq:v85-critical-input}; the potential bound follows by
Fubini from \eqref{eq:v83-fixed-width-input}. Division by $s$ proves
\eqref{eq:v85-band-critical}. When $b_0=\sqrt{8\Lambda_\circ s}$,
the exponents are $\alpha/2=129/256$ and $(1+\alpha)/2=257/256$.
All widths and this outgoing restriction were fixed before the
collision limit. Incoming grazing and backward critical ends require
no deletion in this argument.
\end{proof}

Let $\overline{\mathcal F}_r^{\ge\delta}$ retain only $l\ge\delta$.
Let $\mathcal K_r^{\ge\delta}$ further require $l\le2r$ and
$z(b-)^2<8\Lambda_\circ r$. These are restrictions of the original
surface, not different surface definitions.

\begin{theorem}[Full surface height with an arbitrary positive residence cutoff]
\label{thm:v85-cutoff-surface}
For every fixed $r>0$ sufficiently small and every fixed $\delta>0$,
\begin{align}
 \limsup_m\sup_{R,n,k}m^2
       \|\overline{\mathcal F}_r^{\ge\delta}\|_\infty
                       &\le C_\eta r^{1/32},\label{eq:v85-cutoff-surface}\\
 \limsup_m\sup_{R,n,k}m^2\|\mathcal K_r^{\ge\delta}\|_\infty
                       &\le C_\eta r^{129/256}(1+|\log r|).
                                             \label{eq:v85-cutoff-critical}
\end{align}
The leading constants are independent of $\delta$. The paired incoming
assignment on short noncritical segments, restricted to $l\ge\delta$,
also has ordered height at most $C_\eta r^{1/32}$.
The statements fix $\delta$ before the collision limit; they do not
assert these bounds for $\delta=0$.
\end{theorem}
\begin{proof}
For $l>2r$ use the long-residence persistence, reproduced here with
notation distinct from the arithmetic reference:
\[
 \mathcal R_r(t):=\int_{\{l>2r\}}k_r(t;y)\,d\zeta(y)
       \le\frac2r\int_{t-r}^t\overline{\mathcal F}_{2r}(v)\,dv.
\]
Indeed, if $b-t\in(0,r)$, then $k_{2r}(t-u;y)\ge k_r(t;y)/2$
for $0<u<r$. Its ordered height is $C_\eta r^{1/32}$.

For $\delta\le l\le2r$ choose a finite $N$ with $2^{-N}r<\delta$.
The bands $(2^{-j}r,2^{1-j}r]$, $0\le j\le N$, cover these
segments, and $k_r\le\one_{(a,b)}$. Take the collision limsup
of this finite sum and apply \eqref{eq:v85-band-height}. Its leading
constants sum to at most
$C_\eta r^\sigma/(1-2^{-\sigma})$, independently of $N$.
This proves \eqref{eq:v85-cutoff-surface}. No infinite physical
band sum has been moved through a limsup.

For the critical part put $b_r=\sqrt{8\Lambda_\circ r}$ and
$s_j=2^{-j}r$. Apply the smaller of \eqref{eq:v85-band-height}
and \eqref{eq:v85-band-critical} to each of the finitely many bands.
Using $\min(x,y+z)\le\min(x,y)+z$, their sum is at most
\[
 C_\eta\sum_{j\ge0}\min\{s_j^\sigma,b_r^\alpha\}
       +C_\eta b_r^{\alpha-\eta}\sum_{j\ge0}s_j^{(1+\eta)/2}.
\]
These are sums of numerical upper bounds, not interchanged physical
limits. Splitting the first series where $s_j^\sigma$ crosses
$b_r^\alpha$ bounds it by
$C b_r^\alpha(1+|\log r|)$ for small $r$. The second is
$O_\eta(r^{(1+\alpha)/2})$. This proves
\eqref{eq:v85-cutoff-critical}. Finally apply
Corollary~\ref{cor:v85-audited-transfer} to the selected short
noncritical family. Its outgoing source is bounded by
\eqref{eq:v85-cutoff-surface}; its paired incoming source differs
by at most $Cr$. This proves the last height bound.
\end{proof}

The theorem controls the complete surface at every previously fixed
positive residence cutoff, with a uniform leading constant. In
particular it covers both kinds of short surfaces in the preceding
report. The remaining issue is the complete mass carried by residence
lengths tending to zero before the collision limit, not another
unestimated positive residence band. A separate ordered tail must
still be retained, just as for the fractional volume source.
